\chapter{Maxwell方程组}
\label{chap:Maxwell}
\section{Maxwell位移电流}
在前面的讨论中，我们已经得到了关于电场和磁场的基本定律，分别是：
\begin{alignat*}{3}
    \text{电场Gauss定律}&\,\eqref{eqn:divD}&\qqquad\div\bm D &= \rho,\\
    \text{磁场Gauss定律}&\,\eqref{eqn:divB-macro}&\div\bm B &= 0,\\
    \text{Ampère定律}&\,\eqref{eqn:curlH}&\curl\bm H &= \bm J,\\
    \text{Faraday电磁感应定律}&\,\eqref{eqn:curlE-B}&\curl\bm E &= -\pv{\bm B}t.
\end{alignat*}
%除了Faraday定律外，所有定律都是在稳恒态条件下得到的。
但在\chapref{chap:magnetostatics}推导Ampère定律时，我们设定了稳恒电流条件$\div\bm J = 0$，这在非稳恒态问题中便显然不适用了，而应该修正为连续性方程\eqref{eqn:continuity}：%与连续性方程不一致了。
%Maxwell受Faraday观察结果的启发，发现了方程组中的矛盾。
\[
    \pv\rho t + \div\bm J = \div\biggp{\pv{\bm D}t + \bm J} = 0.
\]
上式利用了电场的Gauss定律，从而Ampère定律也被修正为
\begin{equation}
    \label{eqn:curlH-JD}
    \curl\bm H = \bm J + \pv{\bm D}t.
\end{equation}
Maxwell将附加项$\p\bm D/\p t$称为位移电流(displacement current)。
\section{Maxwell方程组}
修正Ampère定律后，我们便得到了Maxwell方程组：\footnote{事实上Maxwell原稿中的方程组是分量形式，并不简洁，其微分形式是由Heaviside给出的。}
\begin{subequations}
    \begin{align}
        \label{eqn:Maxwell divD}
        \div\bm D &= \rho,\\
        \label{eqn:Maxwell divB}
        \div\bm B &= 0,\\
        \label{eqn:Maxwell curlE}
        \curl\bm E &= -\pv{\bm B}t,\\
        \label{eqn:Maxwell curlH}
        \curl\bm H &= \bm J + \pv{\bm D}t.
    \end{align}
\end{subequations}
将Maxwell方程和Lorentz力
\begin{equation}
    \label{eqn:Lorentz}
    \bm F = q(\bm E + \bm v \times \bm B),
\end{equation}
以及Newton第二定律结合起来，就可以完全描写相互作用 带电粒子和电磁场的经典动力学。
%Maxwell方程组是非常自洽、有非常深刻的物理背景。
\begin{example}{真空中电磁波}{electromagnetic wave}
    在无源的真空中
    \begin{gather*}
        \rho = 0,\enspace\bm J = \bm 0;\\
        \bm D = \varepsilon_0\bm E,\enspace\bm H = \bm B/\mu_0;
    \end{gather*}
    利用式\eqref{eqn:curlcurl}，可得
    \begin{align*}
        \laplacian\bm E &= \grad(\div\bm E) - \curl(\curl\bm E) = 0 + \pp t(\curl\bm B)
        = \mu_0\varepsilon_0\pv[2]{\bm E}t,\\
        \laplacian\bm B &= \grad(\div\bm B) - \curl(\curl\bm B) = 0 - \mu_0\varepsilon_0\pp t(\curl\bm E)
        = \mu_0\varepsilon_0\pv[2]{\bm B}t.
    \end{align*}
    因此真空中电场和磁场的各个分量均服从波动方程(wave equation)：
    \begin{equation}
        \laplacian u = \frac1{c^2}\pv[2]ut,
    \end{equation}
    对于自由空间$\RR^3$，有平面波解，且波速为
    \begin{equation}
        c = \frac1{\sqrt{\mu_0\varepsilon_0}} = \qty{299 792 458}{\metre\per\second},
    \end{equation}
    与真空中光速相同！这意味着光就是电磁波(electromagnetic wave)。Maxwell方程组将电、磁和光作为同一现象的表现形式。
\end{example}

\subsection{标量势和矢量势}

Maxwell方程组是4个联立的一阶偏微分方程组，它们将电场和磁场的各个分量联系起来。在简单情况下这一方程组可以解出来，但是引入势往往更方便。

\begin{definition}
    {标量势和矢量势}{scalar potential and vector potential}
    引入标量势$\Phi$和矢量势$\bm A$，使得电场和磁场可以表示为
    \begin{subequations}
        \begin{align}
            \label{eqn:-nablaPhi-pApt}
            \bm E &= -\grad\Phi - \pv{\bm A}t,\\
            \label{eqn:curlA}
            \bm B &= \curl\bm A.
        \end{align}
    \end{subequations}
\end{definition}

\begin{proof}
    式\eqref{eqn:Maxwell divB}可以写成$\div(\curl\bm A) \equiv 0$的形式，
    继而式\eqref{eqn:Maxwell curlE}变为 
    \[
        \curl\bm E + \pv{\bm B}t = \curl\biggp{\bm E + \pv{\bm A}t} = \bm 0,
    \]
    也可以写成$\curl\grad\Phi \equiv \bm 0$的形式。
\end{proof}

\begin{corollary}
    剩下的式\eqref{eqn:Maxwell divD}和\eqref{eqn:Maxwell curlH}可变形为%(考虑微观情形而非宏观平均)
    \begin{subequations}
        \label{eqn:Phi-A coupled}
        \begin{align}
            \label{eqn:Phi-A coupled a}
            \laplacian\Phi + \pp t(\div\bm A) &= -\frac\rho{\varepsilon_0},\\
            \label{eqn:Phi-A coupled b}
            \laplacian\bm A - \frac1{c^2}\pv[2]{\bm A}t - \grad\biggp{\div\bm A + \frac1{c^2}\pv\Phi t} &= -\mu_0\bm J.
        \end{align}
    \end{subequations}
    下一步是解耦合(decouple)式\eqref{eqn:Phi-A coupled}中的$\Phi$和$\bm A$。
\end{corollary}

\begin{theorem}
    {规范不变性}{gauge invariance}
    电场和磁场在规范变换(gauge transformation)中保持不变：
    \begin{equation}
        (\Phi,\bm A) \mapsto \Bigp{\Phi - \pv\varLambda t,\bm A + \grad\varLambda}.
    \end{equation}
\end{theorem}

\begin{proof}
    规范变换后的$\Phi',\bm A'$得到的电场和磁场为
    \begin{align*}
        \bm E' &= -\grad\biggp{\Phi - \pv\varLambda t} - \pp t(\bm A + \grad\varLambda) = -\grad\Phi - \pv{\bm A}t = \bm E;\\
        \bm B' &= \curl(\bm A + \grad\varLambda) = \curl\bm A = \bm B.
        \qedhere
    \end{align*}
\end{proof}

\begin{corollary}
    总能将任意$(\Phi,\bm A)$规范变换为满足特定约束条件的$(\Phi',\bm A')$。
\end{corollary}

\paragraph{Lorenz规范}

一种直接的想法是令梯度内的项为0
\begin{equation}
    \div\bm A + \frac1{c^2}\pv\Phi t = 0,
\end{equation}
这称为Lorenz规范\footnote{注意区分Lorentz和Lorenz。}，从而Maxwell方程组可以写成
\begin{subequations}
    \begin{align}
        \laplacian\Phi - \frac1{c^2}\pv[2]\Phi t &= -\frac\rho{\varepsilon_0};\\
        \laplacian\bm A - \frac1{c^2}\pv[2]{\bm A}t &= -\mu_0\bm J.
    \end{align}
\end{subequations}
上面两式其实形式是一样的，可定义d'Alembertian%\footnote{其实就是四维赝Riemann度量下的Laplace算子。}
\begin{equation}
    \square := \frac1{c^2}\pp[2]t - \laplacian.
\end{equation}
若规范变换前后均满足Lorenz规范，则
\begin{equation}
    \laplacian\varLambda - \frac1{c^2}\pv[2]\varLambda t = 0.
\end{equation}
\paragraph{Coulomb规范}
另一种规范是\secref{sec:vector potential}提到的Coulomb规范
\begin{equation}
    \label{eq:ed-coulomb-gauge}
    \div\bm A = 0,
\end{equation}
得到的方程为
\begin{subequations}
    \begin{align}
        \label{eqn:Coulomb gauge Phi}
        \laplacian\Phi &= -\frac\rho{\varepsilon_0};\\
        \label{eqn:Coulomb gauge A}
        \laplacian\bm A - \frac1{c^2}\pv[2]{\bm A}t &= -\mu_0\bm J + \frac1{c^2}\pp t\grad\Phi.
    \end{align}
\end{subequations}

式\eqref{eqn:Coulomb gauge Phi}是Poisson方程，可直接得到
\begin{equation}
    \label{eq:ed-coulomb-gauge-potential}
    \Phi(\bm r,t) = \frac1{4\pi\varepsilon_0}\int_V\frac{\rho(\bm r',t)}{\abs{\bm r - \bm r'}}\d V'
 \end{equation}
这说明：标量势$\Phi(\bm r,t)$是由电荷密度$\rho(\bm r',t)$产生的瞬时Coulomb势\footnote{这是其被称为Coulomb规范的原因。}，%二者在不同地方同时改变，
它在这个规范下表现为瞬时函数；但势不是直接可观测量，这不表示电磁信号可以瞬时传播。$-\grad\Phi$与$-\p\bm A/\p t$中的非推迟部分相消后，物理电磁场仍满足因果传播。

对于式\eqref{eqn:Coulomb gauge A}，将电流密度$\bm J$写成纵向(longitudinal,或无旋irrotational)电流$\bm J_{\mathrm{l}}$和横向(transverse，或螺旋solenoidal)电流$\bm J_\transv$之和，满足
\[
    \curl\bm J_{\mathrm{l}} = \bm 0,\quad\div\bm J_\transv = 0.
\]
从而将式\eqref{eqn:Coulomb gauge A}解耦合为(略去推导过程)
\begin{subequations}
    \begin{align}
        \frac1{c^2}\pp t\grad\Phi &= \mu_0\bm J_{\mathrm{l}};\\
        \laplacian\bm A - \frac1{c^2}\pv[2]{\bm A}t &= -\mu_0\bm J_\transv.
    \end{align}
\end{subequations}
这说明：标量势方程的源可以用纵向电流$\bm J_{\mathrm{l}}$表示；而矢量势满足波动方程，在空间中以光速$c$传播，其源可以用横向电流$\bm J_\transv$表示。在Coulomb规范下，横向电场由$-\p\bm A/\p t$给出，纵向电场由$-\grad\Phi$给出；将它们组合为可观测场后，才能讨论传播与因果关系。

\begin{example}{Coulomb规范下电磁波的表达}{em wave in Coulomb gauge}
    在全空间无源、无外加静电场并取无穷远处电势为零时，可取$\Phi = 0$；第二个方程
    \[
        \laplacian\bm A - \frac1{c^2}\pv[2]{\bm A}t = \bm 0 \implies \bm A = \bm A_0\e^{\i(\bm k\cdot\bm r-\omega t)}.
    \]
    从而
    \begin{align*}
        \bm B &= \curl\bm A = \i\bm k \times \bm A;\\
        \bm E &= 0 - \pv{\bm A}t = \i\omega\bm A.
    \end{align*}
\end{example}
\paragraph{磁荷存在吗？}
如果自由磁荷真的存在，Maxwell方程组就要改写
\begin{align*}
    \div\bm D &= \rho_\electric,\\
    \div\bm B &= \rho_\magnetic,\\
    \curl\bm E &= -\bm J_\magnetic - \pv{\bm B}t,\\
    \curl\bm H &= \bm J_\electric + \pv{\bm D}t.
\end{align*}
我们需要新的理论，标量势和矢量势的形式也需要改写。
\footnote{事实上，无论磁荷存在与否，理论物理学家都可以提出许多漂亮的理论，但是决定磁荷存在性仍需要实验物理学家的实验验证。举个例子：Higgs粒子是粒子物理学标准模型(standard model)预言的粒子，也有一些Higgsless的理论
%\footnote{唐老师在课上不点名的提到了一位已经离开工物系的大佬，推测可能是高原宁院士。}
大受欢迎，直到2012年Higgs粒子被发现。}
\subsection{波动方程的Green函数}
Lorenz规范下，$\Phi$和$\bm A$均满足有源波动方程
\begin{equation}
    \label{eqn:DAlembertianPsi=4pif}
    \laplacian\Psi - \frac1{c^2}\pv[2]\Psi t = -4\pi f(\bm r,t).
\end{equation}
%其Green函数
%\biggp{\laplacian-\frac1{c^2}\pp[2]t}G(\bm r,t;\bm r',t')=-4\pi\delta(\bm r-\bm r')\delta(t-t'),
%表示在$\bm r'$处的点源，且仅在$t'$时刻激发。
对$t$进行Fourier变换，%频域里
\[
    (\laplacian + k^2)\hat\Psi(\bm r,\omega) = -4\pi\hat f(\bm r,\omega),
\]
其中$k := \omega/c$。
则频域中的Green函数$\hat G_k(\bm r,\bm r')$满足
\begin{equation}
    \label{eq:ed-helmholtz-green-source}
    (\laplacian + k^2)\hat G_k(\bm r,\bm r') = -4\pi\delta(\bm r - \bm r').
\end{equation}
%不同于时域中的Green函数$G(\bm r,t;\bm r',t')$
$\hat G_k$是关于源点$\bm r'$球对称的，定义$\bm R := \bm r - \bm r'$可得
\[
    \frac1R\dd[2]R(R\hat G_k) + k^2\hat G_k = -4\pi\delta(\bm R),
\]
当$R \neq 0$时，
\[
    \pp[2]R(R\hat G_k) + k^2(R\hat G_k) = 0 \implies \hat G_k^\pm(R) = \frac{\e^{\pm\i kR}}R
\]
%即Green函数可以写成
有两个线性无关的球面波(spherical wave)解，
分别为发散(diverging)球面波$\hat G_k^+(R)$和会聚(converging)球面波$\hat G_k^-(R)$，后者的例子是透镜聚焦。则频域下Green函数是通解的线性组合
\[
    \hat G_k(R) = A\hat G_k^+(R) + B\hat G_k^-(R),
\]
Green函数要求在点源处
\[
    \lim_{kR\to0}\hat G_k(R) = \frac1R \implies A + B = 1.
\]
回到时域下的Green函数，满足
\begin{equation}
    \label{eq:ed-wave-green-source}
    \biggp{\laplacian - \frac1{c^2}\pp[2]t}G^\pm(\bm r,t;\bm r',t') = -4\pi\delta(\bm r - \bm r')\delta(t - t'),
\end{equation}
对$t$进行Fourier变换， 
\[
    (\laplacian + k^2)\hat G(\bm r,\omega;\bm r',t') = -4\pi\delta(\bm r - \bm r')\e^{\i\omega t'}.
\]
可得 
\[
    \hat G^\pm(\bm r,\omega;\bm r',t') = \hat G_k^\pm(R)\e^{\i\omega t'} = \frac{\e^{\i(\pm kR+\omega t')}}R,
\]
进行Fourier反变换，
\[
    G^\pm(R,\tau) = \frac1{2\pi}\intR\frac{\e^{\i(\pm kR-\omega\tau)}}R\d\omega,\quad\tau := t - t',
\]
对于非色散(nondispersive)介质，$c$与频率无关，Green函数有通解
\begin{equation}
    \label{eq:ed-causal-wave-green}
    G^\pm(R,\tau) = \frac1R\delta\biggp{\tau \mp \frac Rc}
\end{equation}
即
\begin{equation}
    G^\pm(\bm r,t;\bm r',t') = \frac1{\abs{\bm r - \bm r'}}\delta\biggp{t' - \Bigp{t \mp \frac{\abs{\bm r - \bm r'}}c}}.
\end{equation}
分别为推迟(retarded)和超前(advanced) Green函数。$\delta$函数的自变量表明：在$\bm r$点于时间$t$观测的效应是由距离为$\bm R$的源在%推迟(或超前)
时间为$t' = t \mp R/c$时的作用引起的。\footnote{课上的理解是：无穷远处产生的源经过$\bm r$后仍需要$R/c$的时间去经过$\bm r'$。}
从而有源波动方程\eqref{eqn:DAlembertianPsi=4pif}的两个解为 
\[
    \Psi^\pm(\bm r,t) = \int G^\pm(\bm r,t;\bm r',t')f(\bm r',t')\d V'\nd t',
\]
我们考虑一个在时间和空间中作局域分布的源$f(\bm r',t')$，这个源分布只在$t' = 0$前后一段有限的时间间隔内不为零。则在无源的波动方程仅仅满足其次波动方程
\[
    \laplacian\Psi - \frac1{c^2}\pv[2]\Psi t = 0,
\]
想像两种极限情形：
\begin{compactenum}
	\item $t \to -\infty$，存在一个特定的“背景”波$\Psi_\text{in}$，
	\begin{equation}
	    \label{eq:ed-incoming-retarded-decomposition}
        \Psi(\bm r,t) = \Psi_\text{in}(\bm r,t) + \int G^+(\bm r,t;\bm r',t')f(\bm r',t')\d V'\nd t';
    \end{equation}
	\item $t \to +\infty$，存在一个特定的“稳态”波$\Psi_\text{out}$，
	\begin{equation}
	    \label{eq:ed-outgoing-advanced-decomposition}
        \Psi(\bm r,t) = \Psi_\text{out}(\bm r,t) + \int G^-(\bm r,t;\bm r',t')f(\bm r',t')\d V'\nd t',
    \end{equation}
\end{compactenum}
最常见的物理情景是$\Psi_\text{in} = 0$%“超前”这个概念总是太科幻了，
\begin{subequations}
    \begin{align}
        \Psi(\bm r,t) &= \int\frac{f(\bm r',t')}{\abs{\bm r - \bm r'}}\d V',\\
        t' &= t - \frac{\abs{\bm r - \bm r'}}c.
    \end{align}
\end{subequations}
物理学家已经就一维、二维和三维情形下的有限时间的初值和终值问题进行过广泛研究。\footnote{Morse and Feshbach, Methods of Theoretical Physics, pp.~843--847.}

\subsection{推迟势和推迟场}


如果电荷密度和电流密度分布随时间变化，观测点的标量势和矢量势将会延迟：
\begin{subequations}
    \begin{align}
        \Phi(\bm r,t) &= \frac1{4\pi\varepsilon_0}\int\frac{\rho(\bm r',t')}{\abs{\bm r - \bm r'}}\d V',\\
        \bm A(\bm r,t) &= \frac{\mu_0}{4\pi}\int\frac{\bm J(\bm r',t')}{\abs{\bm r - \bm r'}}\d V'.
    \end{align}
\end{subequations}
% 由
% 但电场和磁场并不能简单的变换\eqref{eqn:E(x)}和\eqref{eqn:B(x)}中的$\rho,\bm J$，而需要从
% \[
%     \bm E=-\grad\Phi-\pv{\bm A}t,\quad\bm B=\curl\bm A
% \]
得到
\begin{subequations}
    \begin{align}
        \bm E(\bm r,t) &= \frac1{4\pi\varepsilon_0}\int\biggb{-[\nabla'\rho]_\ret - \frac1{c^2}\pv{\bm J}{t'}}\frac{\d V'}{\abs{\bm r - \bm r'}},\\
        \bm B(\bm r,t) &= \frac{\mu_0}{4\pi}\int\frac{[\nabla' \times \bm J]_\ret}{\abs{\bm r - \bm r'}}\d V',
    \end{align}
\end{subequations}
% 注意上式中$\nabla'\rho$和$\nabla'\times\bm J$是先求梯度和散度再令$t'=t-R/c$。
Jackson用$[\cdot]_\ret$表示赋值$t' = t - R/c$，这样的好处是可以有效区分$\grad[f]_\ret \neq [\grad f]_\ret$，因为前者的$t$与$\bm r'$有关，进一步的：
\begin{align*}
    [\nabla'\rho]_\ret &= \nabla'[\rho]_\ret - \Bigb{\pv\rho{t'}}_\ret\nabla'(t - R/c);\\
    [\nabla' \times \bm J]_\ret &= \nabla' \times [\bm J]_\ret + \Bigb{\pv{\bm J}{t'}}_\ret \times \nabla'(t - R/c).
\end{align*}
而$\nabla'R = -{\uvec R}$，故
\begin{subequations}
    \begin{align}
        \bm E(\bm r,t) &= \frac1{4\pi\varepsilon_0}\int\biggb{\biggp{\frac{\rho(\bm r',t')}R + \frac1c\pv\rho{t'}}\uvec R - \frac1{c^2}\pv{\bm J}{t'}}\frac{\d V'}R,\\
        \bm B(\bm r,t) &= \frac{\mu_0}{4\pi}\int\biggb{\frac{\bm J(\bm r',t')}R + \frac1c\pv{\bm J}{t'}} \times \frac{\uvec R}R\d V'.
    \end{align}
\end{subequations}
上式称为Jefimenko对Coulomb定律和Biot--Savart定律的推广。
若$\rho$和$\bm J$是与时间无关的，则可重新得到静电场和静磁场的表达式\eqref{eqn:E(r)}和\eqref{eqn:B(r)}。
\begin{example}{电荷量变化的点电荷}{time-varying point charge}
    设电流流入或流出原点的小电极，使点电荷项随时间变化。其电荷密度为
    \[
        \rho(\bm r',t') = Q(t')\delta(\bm r'),
    \]
    由连续性方程，供电电流的散度在原点有一项$-\dot Q(t)\delta(\bm r)$。该点电荷对Lorenz规范推迟标势的贡献为
    \[
        \Phi(r,t) = \frac1{4\pi\varepsilon_0r}Q\Bigp{t - \frac rc}.
    \]
    求完整电场时，还须计入其余电荷的标势以及供电电流的矢量势。
\end{example}
\begin{example}{阶跃电流}{current}
    无限长直导线通阶跃电流
    \[
        I(t) = \begin{cases}
            0,&t \leqslant 0\\
            I_0,&t > 0
        \end{cases}
    \]
    电流密度为
    \[
        \bm J(\bm r',t') = I(t')\delta(x')\delta(y')\uvec z.
    \]
    矢量势
    \begin{align*}
        \bm A(\bm r,t) &= \frac{\mu_0}{4\pi}\uvec z\intR I\Bigp{t - \frac Rc}\frac{\d z'}R,\\
        R &= \sqrt{\rho^2 + (z - z')^2},
    \end{align*}
    $t \leqslant \rho/c$时，$\bm A = \bm 0$；
    $t > \rho/c$时，仅$(z' - z)^2 \leqslant (ct)^2 - \rho^2$部分有贡献，则
    \[
        \bm A(\rho,t) = \frac{\mu_0I_0}{2\pi}\ln\biggb{\frac{ct + \sqrt{(ct)^2 - \rho^2}}{\rho}}\uvec z.
    \]
    由式\eqref{eqn:-nablaPhi-pApt}和\eqref{eqn:curlA}，电场和磁场分别为
    \begin{subequations}
        \label{eq:ed-switched-wire-fields}
        \begin{align}
            \bm E(\rho,t) &= -\frac{\mu_0I_0}{2\pi}\frac c{\sqrt{(ct)^2 - \rho^2}}\uvec z,\\
            \bm B(\rho,t) &= \frac{\mu_0I_0}{2\pi}\frac{ct}{\rho\sqrt{(ct)^2 - \rho^2}}\uvec\phi.
        \end{align}
    \end{subequations}
    $t \to +\infty$时，电磁场趋于稳定，与恒稳电流情形相同：
    \[
        \bm E = \bm 0,\quad\bm B = \frac{\mu_0I_0}{2\pi\rho}\uvec\phi.
    \]
\end{example}
\subsection{电磁Green张量与互易定理}
\label{sec:dyadic-green}
给定一个振荡电偶极，在另一点放置一根短探针，探针沿哪个方向才能测到电场？前面标量波动方程的Green函数给出了推迟时间。要直接从偶极矩求矢量电场，还须保留源与探针的方向。取时间因子$\e^{-\i\omega t}$，记$k = \omega/c$。在非磁、线性、空间局域的介质中，相对介电张量为$\bm\varepsilon_\relative(\bm r,\omega)$，Maxwell方程给出
\begin{equation}
    \label{eq:ed-electric-wave-operator}
    \curl(\curl\bm E) - k^2\bm\varepsilon_\relative \cdot \bm E = \i\omega\mu_0\bm J .
\end{equation}
记左边的算符为$\mathcal L$。在给定边界条件和向外辐射条件下，定义电场的推迟Green张量
\begin{subequations}\label{eqn:electric dyadic green}
    \begin{align}
        \mathcal L_{\bm r}\bm G(\bm r,\bm r';\omega) &=
        \bm I\,\delta^{(3)}(\bm r - \bm r'),\\
        \bm E(\bm r) &= \i\mu_0\omega\int\bm G(\bm r,\bm r';\omega) \cdot \bm J(\bm r')\,\d V'.
    \end{align}
\end{subequations}
这里$\bm I$为单位张量，$\bm J$为给定的外加电流密度。位于$\bm r_0$的点偶极$\bm p$产生电流$\bm J = -\i\omega\bm p\,\delta^{(3)}(\bm r - \bm r_0)$，故探针测得的第$i$个分量是$\textstyle E_i = \mu_0\omega^2\sum_j G_{ij}p_j$。

真空中取$R = \abs{\bm r - \bm r_0}$、$\uvec n = (\bm r - \bm r_0)/R$。令$g = \e^{\i kR}/(4\pi R)$，它满足$(\laplacian + k^2)g = -\delta^{(3)}(\bm r - \bm r_0)$。利用$\curl(\curl\bm E) = \grad(\div\bm E) - \laplacian\bm E$，可验得$\bm G_0 = (\bm I + k^{-2}\grad\grad)g$，两个梯度都对场点求导。于是在$R > 0$处，
\begin{equation}
    \bm E(\bm r) = \frac{\e^{\i kR}}{4\pi\varepsilon_0}
    \biggl\{\frac{k^2}{R}\bigb{\bm p - \uvec n(\uvec n \cdot \bm p)}
    + \biggl(\frac{1}{R^3} - \frac{\i k}{R^2}\biggr)
    \bigb{3\uvec n(\uvec n \cdot \bm p) - \bm p}\biggr\}.
    \label{eqn:dyadic dipole field}
\end{equation}
近区的主项恢复静电偶极场。对于线性振荡的偶极，沿偶极轴观测，$1/R$项为零。在赤道面观测，这一项的电场沿偶极方向。

\paragraph{交换源和探针}
令$\bm J_1,\bm J_2$为两个同频电流源，用相同下标标记各自的场。将两组Maxwell方程交叉点乘并相减，若$\bm\varepsilon_\relative^\top = \bm\varepsilon_\relative$，得到
\[
    \div(\bm E_1 \times \bm H_2 - \bm E_2 \times \bm H_1)
      = \bm E_2 \cdot \bm J_1 - \bm E_1 \cdot \bm J_2 .
\]
理想导体壁上的切向电场为零，远处两组外行波的表面项相消。体积分后有
\begin{equation}
    \label{eq:ed-lorentz-reciprocity-sources}
    \int\bm E_2 \cdot \bm J_1\,\d V = \int\bm E_1 \cdot \bm J_2\,\d V.
\end{equation}
用两个取向任意的点偶极作源，得到Lorentz互易定理
\begin{equation}
    G_{ij}(\bm r,\bm r';\omega) = G_{ji}(\bm r',\bm r;\omega).
\end{equation}
源和探针互换时要同时交换两者的取向。推导中的点积没有复共轭，故具有对称复介电张量的吸收介质也可满足互易性。

\paragraph{理想镜面旁的偶极}
理想导体占据$z < 0$，真空中的偶极位于$\bm r_0 = (0,0,h)$。给定偶极矩的峰值复振幅$\bm p$，驱动源维持振荡所付出的时均功率为
\[
    P = -\frac12\Re\int\bm J\conj \cdot \bm E\,\d V
      = \frac{\omega}{2}\Im\bigb{\bm p\conj \cdot \bm E(\bm r_0)} .
\]
导体无损，向上半空间辐射的功率便等于$P$。把$\bm G$分成真空项$\bm G_0$和反射项$\bm G_\mathrm{ref}$。重合点处真空Green张量的实部有奇性，计算功率所用的虚部却是有限的。展开$\sin(kR)/(4\pi R)$并对它作用$\bm I + k^{-2}\grad\grad$，得$\Im\bm G_0(\bm r_0,\bm r_0) = k\bm I/(6\pi)$。因此
\begin{subequations}
    \label{eq:ed-dipole-environment-power}
    \begin{align}
        P_0 &= \frac{\omega^4\abs{\bm p}^2}{12\pi\varepsilon_0 c^3},\\
        P - P_0 &= \frac{\mu_0\omega^3}{2}
           \Im\bigb{\bm p\conj \cdot \bm G_\mathrm{ref}(\bm r_0,\bm r_0) \cdot \bm p}.
    \end{align}
\end{subequations}
镜像偶极位于$(0,0,-h)$，其切向分量反号、法向分量同号，使导体上的总切向电场为零。在式\eqref{eqn:dyadic dipole field}中放入镜像偶极，取$R = 2h$、$\uvec n = \uvec z$，便得到作用回源点的反射场。记往返相位$x = 2kh$，由该场对源做的功可得
\begin{equation}
    \begin{aligned}
        \frac{P_\perp}{P_0} &= 1 + 3\biggl(\frac{\sin x}{x^3} - \frac{\cos x}{x^2}\biggr),\\
        \frac{P_\parallel}{P_0} &= 1 - \frac32\biggl(
            \frac{\sin x}{x} + \frac{\cos x}{x^2} - \frac{\sin x}{x^3}\biggr).
    \end{aligned}
    \label{eqn:mirror-dipole-rates}
\end{equation}
图\ref{fig:mirror-dipole-power}中的下标以镜面为准。靠近镜面时，平行偶极和镜像在远区相消，$P_\parallel/P_0 \simeq x^2/5$。垂直偶极与镜像同相，在上半空间场振幅加倍，积分所得功率趋于$2P_0$。远离镜面时，两种取向的总功率都趋于$P_0$。\chapref{chap:cavity qed}将用相同的两个比值计算原子在镜前的自发发射寿命。\footnote{电磁Green张量与自发发射率的关系可参阅S. Scheel、L. Knöll和D.-G. Welsch，\href{https://arxiv.org/abs/quant-ph/9905007}{Spontaneous decay in the presence of absorbing dielectric bodies} (1999)。}
\begin{figure}[htbp]
    \centering
    \includegraphics[width=.71\linewidth]{figures/mirror_dipole_power.pdf}
    \caption{保持偶极矩振幅不变时，理想镜面前两种取向的总辐射功率。$h$为偶极与镜面的距离，$\lambda$为真空波长。}
    \label{fig:mirror-dipole-power}
\end{figure}

\section{宏观电磁场}
尽管我们已经得到了Maxwell方程组、电磁场与标量势矢量势的关系，但我们仍没有处理$\bm B\&\bm H,\bm D\&\bm E$之间的关系。在真空中
\begin{gather*}
    \bm D = \varepsilon_0\bm E,\enspace\bm B = \mu_0\bm H,\\
    c^2 = \frac1{\varepsilon_0\mu_0}.
\end{gather*}
有质介质中，
\begin{equation}
    \label{eq:ed-macroscopic-constitutive-definitions}
    \bm D = \varepsilon_0\bm E + \bm P,\quad \bm H = \bm B/\mu_0 - \bm M.
\end{equation}
对于线性介质，$\bm P = \varepsilon_0\chi_\electric\bm E$，
\begin{equation}
    \label{eq:ed-isotropic-constitutive-relations}
    \bm D = \varepsilon\bm E,\quad\bm H = \bm B/\mu;
\end{equation}
在非线性介质中，关系变得复杂。
\paragraph{微观场}
考虑一个仅由电子和原子核构成的微观世界，由于原子之内是如此的空旷，以至于可以认为是真空条件，因此微观场的Maxwell方程
\begin{align*}
    \div\bm e &= \frac\varrho{\varepsilon_0},\\
    \div\bm b &= 0,\\
    \curl\bm e + \pv{\bm b}t &= \bm 0,\\
    \curl\bm b - \frac1{c^2}\pv{\bm e}t &= \mu_0\bm j.
\end{align*}
$\varrho,\bm j$是微观电荷密度和电流密度，$\bm e,\bm b$是微观电场和磁场，而$\bm d,\bm h$并无定义。
\paragraph{宏观场}
宏观场被定义为微观场的空间平均：
\[
    \bm E(\bm r,t) = \ave{\bm e(\bm r,t)},\quad \bm B(\bm r,t) = \ave{\bm b(\bm r,t)},
\]
其中空间平均值是相对于测试函数$f(\bm r)$而言的，其定义为：
\begin{equation}
    \label{eq:ed-spatial-averaging}
    \ave{F(\bm r,t)} := \int F(\bm r - \bm r',t)f(\bm r')\d V',
\end{equation}
因此，空间和时间的微分运算与平均运算可交换：
\[
    \pp{x_\alpha}\ave{F} = \biggave{\pv F{x_\alpha}},\quad\pp t\ave{F} = \biggave{\pv Ft}.
\]
从而宏观场的Maxwell方程为： 
\begin{align}
    \label{eqn:average rho}
    \varepsilon_0\div\bm E &= \ave{\varrho},\\
    \notag
    \div\bm B &= 0,\\
    \notag
    \curl\bm E + \pv{\bm B}t &= \bm0,\\
    \label{eqn:average j}
    \frac1{\mu_0}\curl\bm B - \varepsilon_0\pv{\bm E}t &= \ave{\bm j}.
\end{align}
因此导出场$\bm D$和$\bm H$是
通过从$\ave\varrho$和$\ave{\bm j}$提出某些可以看作介质的宏观性质的贡献而引进的。
下面便研究$\ave\varrho$和$\ave{\bm j}$。
\paragraph{电荷密度}
我们将电荷$\varrho$分为自由电荷$\varrho_{\mathrm{free}}$和束缚电荷$\varrho_{\mathrm{bound}}$，后者被束缚在分子中，因此
\begin{align*}
    \varrho_{\mathrm{free}}(\bm r,t) &= \sum_jq_j\delta(\bm r - \bm r_j);\\
    \varrho_{\mathrm{bound}}(\bm r,t) &= \sum_n\varrho_n(\bm r,t).
\end{align*}
$\varrho_n$是第$n$个分子的电荷密度，分子质心坐标为$\bm r_n$，分子电荷相对质心坐标为$\bm r_{nj}$，则其平均值为
\[
    \ave{\varrho_n(\bm r,t)} = \sum_{j(n)}q_jf(\bm r - \bm r_n - \bm r_{nj}),
\]
因为$\bm r_{nj}$具有原子线度的数量级，适合在$(\bm r - \bm r_n)$处Taylor展开：
\[
    f(\bm r - \bm r_n - \bm r_{nj}) = f(\bm r - \bm r_n) - \bm r_{nj} \cdot \grad f(\bm r - \bm r_n) + \frac12\sum_{\alpha,\beta}(\bm r_{nj})_\alpha(\bm r_{nj})_\beta\pw f{x_\alpha}{x_\beta}(\bm r - \bm r_n) + \cdots,
\]
式中各项分别对应分子电荷、偶极矩、四极矩……
\begin{align*}
    \ave{\varrho_n(\bm r,t)} = {}&q_nf(\bm r - \bm r_n) - \bm p_n \cdot \grad f(\bm r - \bm r_n) + \frac16\sum_{\alpha,\beta}(Q_n)_{\alpha\beta}\pw f{x_\alpha}{x_\beta}(\bm r - \bm r_n) + \cdots
\end{align*}
将上式各项表示成空间平均的形式：
\begin{align*}
    \ave{\varrho_n(\bm r,t)} = {}&\ave{q_n\delta(\bm r - \bm r_n)} - \div\ave{\bm p_n\delta(\bm r - \bm r_n)} + \frac16\sum_{\alpha,\beta}\pw{}{x_\alpha}{x_\beta}\ave{(Q_n)_{\alpha\beta}\delta(\bm r - \bm r_n)} + \cdots
\end{align*}
就平均而言，可以把分子看作位于分子质心(也可选取其他固定点)上的点多极子的集合。对所有分子求和，得到
\[
    \ave{\varrho(\bm r,t)} = \rho(\bm r,t) - \div\bm P(\bm r,t) + \sum_{\alpha,\beta}\pw{}{x_\alpha}{x_\beta}Q_{\alpha\beta}(\bm r,t) + \cdots
\]
其中$\rho,\bm P,Q_{\alpha\beta}$分别是宏观电荷密度、宏观电极化强度、宏观电四极矩密度……
因此式\eqref{eqn:average rho}可以写成 
\[
    \div\bm D = \rho,
\]
宏观电位移矢量定义为
\begin{equation}
    \label{eq:ed-displacement-quadrupole}
    D_\alpha := \varepsilon_0E_\alpha + P_\alpha - \sum_\beta\pv{Q_{\alpha\beta}}{x_\beta} + \cdots
\end{equation}
前两项是熟知的结果，更高阶项原则上是存在的，但几乎可以忽略不计。
\paragraph{电流密度}
$\ave{\bm j}$的推导比$\ave{\varrho}$的处理复杂的多，在这里仅给出结果。定义分子磁矩和宏观磁化强度分别为
\begin{align*}
    \bm m_n &= \frac12\sum_{j(n)}q_j\bm r_{jn} \times \bm v_{jn},\\
    \bm M(\bm r,t) &= \sum_n\bm m_n\delta(\bm r - \bm r_n),
\end{align*}
宏观电流密度
\begin{align*}
    \bm J(\bm r,t) &= \ave{\bm j_{\mathrm{free}} + \bm j_{\mathrm{bound}}},\\
    \bm j_{\mathrm{free}}(\bm r,t) &:= \sum_jq_j\bm v_j\delta(\bm r - \bm r_j).
\end{align*}
进而定义磁场强度\footnote{本式其实是运动介质的%Minkowski电动力学的
非相对论性极限。}
\[
    \bm H := \frac{\bm B}{\mu_0} - \bm M - (\bm D - \varepsilon_0\bm E) \times \bm v.
\]
因此式\eqref{eqn:average j}可以写成
\[
    \curl\bm H - \pv{\bm D}t = \bm J.
\] 

\section{电磁场的能量、动量和角动量}

\paragraph{能量守恒}

电流从场中获得的功率密度为$\bm J \cdot \bm E$。先用宏观场描述静止介质，$\bm J$为显式列入Ampère定律的自由电流。利用Maxwell方程组和式\eqref{eqn:divAxB}得
\[
    \div(\bm E \times \bm H) + \bm E \cdot \pv{\bm D}t
       + \bm H \cdot \pv{\bm B}t = -\bm J \cdot \bm E .
\]
这一步只用到了场方程。进一步设介质的电、磁响应线性、互易、无色散，本构参数为与时间无关的实数或实对称张量。此时上式中两个含时间导数的项恰好合成$\p u/\p t$，其中
\begin{subequations}
    \label{eq:ed-macroscopic-energy-and-flux}
    \begin{align}
        u &= w_\electric + w_\magnetic
          = \frac12(\bm E \cdot \bm D + \bm H \cdot \bm B),\\
        \bm S &= \bm E \times \bm H .
    \end{align}
\end{subequations}
\begin{theorem}{Poynting定律(电磁场的能量守恒定律)}{Poynting theorem}
    体积$V$内的储能增加率，加上通过边界向外流出的功率，等于场对该体积内电流做功的功率取负号。取面积元指向外侧，
    \begin{subequations}
        \label{eq:ed-poynting-theorem}
        \begin{align}
            \pv ut + \div\bm S &= -\bm J \cdot \bm E,\\
            \dd t\int_V u\d V + \oint_{\p V}\bm S \cdot \d\bm A
                 &= -\int_V\bm J \cdot \bm E\d V .
        \end{align}
    \end{subequations}
\end{theorem}
若改用微观场，所有电荷的运动都计入总电流$\bm J_{\mathrm{tot}}$。直接从微观Maxwell方程得到与式\eqref{eq:ed-poynting-theorem}相同形式的守恒律，此时
\begin{subequations}
    \label{eq:ed-microscopic-energy-and-flux}
    \begin{align}
        u_{\mathrm f} &= \frac{\varepsilon_0}2 E^2 + \frac{B^2}{2\mu_0},\\
        \bm S_{\mathrm f} &= \frac{\bm E \times \bm B}{\mu_0}.
    \end{align}
\end{subequations}
即使体积内有电荷，这两个场量仍然适用。若粒子不穿过边界且只有电磁力做功，$\textstyle \int_V\bm J_{\mathrm{tot}} \cdot \bm E\d V$正好是粒子机械能的增加率。宏观描述则把可逆极化和磁化所储存的能量计入$u$，显式传导电流的Joule热计入$\bm J \cdot \bm E$。
\begin{example}{导线的Joule热}{Joule heating of a wire}
    圆柱形导线长度为$L$，半径$a$，两端电压为$V$，则电场强度
    \[
        E = \frac VL,
    \]
    导体表面磁场
    \[
        B = \frac{\mu_0}{2\pi}\frac Ia.
    \]
    Poynting矢量向里，大小为
    \[
        S = \frac VL \cdot \frac1{\mu_0}\frac{\mu_0}{2\pi}\frac Ia = \frac{VI}{2\pi aL},
    \]
    取导线表面的外法向，流入导线并转化为热的功率为
    \[
        -\oint_{\p V}\bm S \cdot \d\bm A = S \cdot 2\pi aL = VI.
    \]
\end{example}
\paragraph{动量守恒}
以下使用真空中的微观场及总电荷、电流密度，$\bm H = \bm B/\mu_0$。所得场动量密度不应直接推广为一般物质中场与物质动量的划分。
电磁场对带电粒子的力是Lorentz力
\[
    \bm F = q(\bm E + \bm v \times \bm B).
\]
则粒子动量的变化率为
\[
    \dv{\bm P_{\mathrm{m}}}t = \int_V(\rho\bm E + \bm J \times \bm B)\d V
\]
消去源，得到
\begin{align*}
    \rho\bm E + \bm J \times \bm B &= \varepsilon_0(\div\bm E)\bm E + \biggp{\frac1{\mu_0}\curl\bm B - \varepsilon_0\pv{\bm E}t} \times \bm B\\
    &= \varepsilon_0\Bigb{\bm E(\div\bm E) - c^2\bm B \times (\curl\bm B) - \bm E \times (\curl\bm E) - \pp t(\bm E \times \bm B)},
\end{align*}
右边出现了%Poynting矢量对
时间的导数，我们依此定义电磁场的动量
\begin{equation}
    \label{eq:ed-vacuum-field-momentum}
    \bm P_{\mathrm{f}} := \varepsilon_0\int_V\bm E \times \bm B\d V = \frac1{c^2}\int_V\bm E \times \bm H\d V.
\end{equation}
电磁场的动量密度
\begin{equation}
    \label{eq:ed-vacuum-momentum-density}
    \bm g := \frac1{c^2}\bm E \times \bm H = \frac1{c^2}\bm S.
\end{equation}
则
\[
    \dv{\bm P_{\mathrm{m}}}t + \dv{\bm P_{\mathrm{f}}}t = \varepsilon_0\int_V\Bigb{\bm E(\div\bm E) - \bm E \times (\curl\bm E) - c^2\bm B \times (\curl\bm B)}\d V.
\]
右边的体积分完全可以加上$c^2\bm B(\div\bm B) \equiv 0$以更对称。并且注意到
\[
    [\bm E(\div\bm E) - \bm E \times (\curl\bm E)]_\alpha = \sum_\beta\pp{x_\beta}\Bigp{E_\alpha E_\beta - \frac12E^2\delta_{\alpha\beta}};
\]
从而右边的体积分具有散度的形式，定义Maxwell应力张量(Maxwell stress tensor)
\begin{equation}
    \label{eqn:vacuum maxwell stress}
    T_{\alpha\beta} := \varepsilon_0\Bigb{E_\alpha E_\beta + c^2B_\alpha B_\beta - \frac12(E^2 + c^2B^2)\delta_{\alpha\beta}}.
\end{equation}
则
\begin{equation}
    \label{eq:ed-integral-momentum-balance}
    \dd t({\bm P_{\mathrm{m}}} + {\bm P_{\mathrm{f}}})_\alpha = \oint_{\p V}\sum_\beta T_{\alpha\beta}\hat n_\beta\d A.
\end{equation}
\paragraph{角动量守恒}
取固定原点，记Lorentz力密度为$\bm f = \rho\bm E + \bm J \times \bm B$。动量守恒的局域形式为
\begin{equation}
    \label{eq:ed-local-momentum-balance}
    \p_t g_i + f_i = \sum_j\p_j T_{ij}.
\end{equation}
场的角动量密度为$\bm r \times \bm g$，力矩密度为$\bm r \times \bm f$。将动量方程对原点取力矩，由于$T_{ij} = T_{ji}$，对$\bm r$求导产生的项$\textstyle \sum_{jk}\epsilon_{ijk}T_{kj}$为零。对固定体积$V$积分得
\begin{equation}
    \dd t\int_V\bm r \times \bm g\,\d V
    + \int_V\bm r \times \bm f\,\d V
    = \oint_{\p V}\bm r \times (\bm T \cdot \uvec n)\,\d A .
\end{equation}
这里$\uvec n$为外法向。按本节的应力符号，$-\bm r \times (\bm T \cdot \uvec n)$为单位面积上向外的角动量通量。对封闭面积分给出电磁场通过边界带走的角动量率。

\begin{example}{Feynman圆盘}{paradox of Feynman disc}
    绝缘圆盘的边缘均匀带有总电荷$Q$，可以绕轴自由转动。记同轴线圈穿过整个盘面的磁通量为$\varPhi(t)$，取$+z$为磁通和转动的正向，忽略圆盘转动产生的自场。电流缓慢关断时，Faraday感应电场给盘缘的电荷以力矩。设盘缘半径为$b$，由轴对称性
    \begin{align*}
        2\pi b E_\phi(b) &= -\dv{\varPhi}t,\\
        N_z &= QbE_\phi(b) = -\frac{Q}{2\pi}\dv{\varPhi}t .
    \end{align*}
    若初始磁通量为$\varPhi_0$，最终为零，则盘得到的角动量为$\Delta L_z = Q\varPhi_0/(2\pi)$。初始的静电场和静磁场已经携带角动量。改变磁通量时，感应电场可以把角动量转移给圆盘。实际装置的完整角动量收支还包括线圈、支架和穿过所选边界的场。
\end{example}
\subsection{色散介质中的Poynting定律}
\label{ssec:Poynting in dissipative media}
在有色散介质中，$\varepsilon,\mu$与频率$\omega$有关，前面的假设不再成立，但仍有
\[
    \int_V\bm J \cdot \bm E\d V = -\int_V\Bigb{\div(\bm E \times \bm H) + \bm H \cdot \pv{\bm B}t + \bm E \cdot \pv{\bm D}t}\d V.
\]
为了讨论色散，有必要利用Fourier变换将时域变到频域中：
\begin{align*}
    \bm E(\bm r,t) &= \intR\hatbm E(\bm r,\omega)\e^{-\i\omega t}\d\omega,\\
    \bm D(\bm r,t) &= \intR\hatbm D(\bm r,\omega)\e^{-\i\omega t}\d\omega;
\end{align*}
在此我们仍假设介质是各向同性的，从而
\[
    \hatbm D(\bm r,\omega) = \varepsilon(\omega)\hatbm E(\bm r,\omega),
\]
其中$\varepsilon$是复数，而场$\bm E(\bm r,t)$必然仍是实数，因此%简单的将积分变量$\omega$变为$-\omega$可得
\[
    \hatbm E(\bm r,-\omega) = \hatbm E\conj(\bm r,\omega),\quad\hatbm D(\bm r,-\omega) = \hatbm D\conj(\bm r,\omega),
\]
从而$\varepsilon(-\omega) = \varepsilon\conj(\omega)$。频率响应使同一点的$\bm D(t)$依赖于此前的电场。即使介质透明，$\bm E \cdot \p\bm D/\p t$通常也不能写成$\p(\bm E \cdot \bm D/2)/\p t$。

\paragraph{窄带波包的储能和吸收}
取中心频率为$\omega_0$的窄带波包，记物理电场为$\Re[\bm E_{\mathrm{a}}(t)\e^{-\i\omega_0t}]$，其中$\bm E_{\mathrm{a}}$是缓变的峰值复振幅，磁场同理。将$-\i\omega\varepsilon(\omega)$在$\omega_0$处展开到一阶，$\p\bm D/\p t$的复包络为
\[
    -\i\omega_0\varepsilon(\omega_0)\bm E_{\mathrm{a}}
        + \eval{\dv{(\omega\varepsilon)}\omega}_{\omega_0}\pv{\bm E_{\mathrm{a}}}t .
\]
在一个载波周期上取平均，相乘的两个实场给出各自复振幅乘积的一半实部。若频带透明，或吸收很弱且吸收项在带宽内的变化可忽略，则
\[
    \biggave{\bm E \cdot \pv{\bm D}t} \simeq
      \frac{\omega_0}{2}\Im\varepsilon\,\abs{\bm E_{\mathrm{a}}}^2 +
      \frac14\dv{(\omega\Re\varepsilon)}\omega\,\pv{\abs{\bm E_{\mathrm{a}}}^2}t .
\]
以下本构参数及频率导数均取在$\omega_0$。加上磁场项，得到包含材料可逆响应的周期平均储能密度（Brillouin公式）和介质吸收的平均功率密度，
\begin{align}
    \label{eqn:Brillouin energy}
    u_\eff &= \frac14\biggb{\dv{(\omega\Re\varepsilon)}\omega\abs{\bm E_{\mathrm{a}}}^2 +
                              \dv{(\omega\Re\mu)}\omega\abs{\bm H_{\mathrm{a}}}^2},\\
    q_{\mathrm{abs}} &= \frac{\omega_0}{2}
        \bigp{\Im\varepsilon\,\abs{\bm E_{\mathrm{a}}}^2 + \Im\mu\,\abs{\bm H_{\mathrm{a}}}^2}.
\end{align}
令$\ave{\bm S} = \Re(\bm E_{\mathrm{a}} \times \bm H_{\mathrm{a}}\conj)/2$，周期平均的能量收支为
\begin{equation}
    \label{eqn:Poynting with loss}
    \pv{u_\eff}t + \div\ave{\bm S} \simeq -\ave{\bm J \cdot \bm E} - q_{\mathrm{abs}} .
\end{equation}
例如，若$\bm J$显式包含$\sigma\bm E$，则$\varepsilon$只计入其余极化响应$\varepsilon_{\mathrm b}$；若改用$\varepsilon = \varepsilon_{\mathrm b} + \i\sigma/\omega$，同一传导损耗已计入$q_{\mathrm{abs}}$，右端的$\bm J$只保留其余显式电流。强吸收频带内，确定可回收的储能还需材料的动力学模型。\footnote{透明频带中Brillouin储能与能量速度的关系见 A.~Welters, Y.~Avniel and S.~G. Johnson, \href{https://doi.org/10.1103/PhysRevA.90.023847}{\textit{Physical Review A} \textbf{90}, 023847 (2014)}, Sec.~II。}

\subsectionstar{谐振场的Poynting定律}
我们假定所有的场和源都是随时间简谐的，即 
\[
    \bm E(\bm r,t) = \Re\bigb{\bm E(\bm r)\e^{-\i\omega t}},
\]
其中$\bm E(\bm r)$一般是复的，有相位的信息。
%因此
%\bm J\cdot\bm E=\frac12\Re\bigb{\bm J\conj(\bm r)\cdot\bm E(\bm r)+\bm J(\bm r)\cdot\bm E(\bm r)\e^{-\i 2\omega t}},
谐振场的Maxwell方程组可以写成
\begin{subequations}
    \begin{align}
        \div\bm D &= \rho,\\
        \div\bm B &= 0,\\
        \label{eqn:Maxwell curlE freq}
        \curl\bm E &= \i\omega\bm B,\\
        \label{eqn:Maxwell curlH freq}
        \curl\bm H &= \bm J - \i\omega\bm D.
    \end{align}
\end{subequations}
取复功率的形式，其实部为场对电流做功的周期平均功率，
\begin{align*}
    \frac12\int_V\bm J\conj \cdot \bm E\d V &= \frac12\int_V(\curl\bm H\conj - \i\omega\bm D\conj) \cdot \bm E\d V\\
    &= \frac12\int_V\Bigb{-\div(\bm E \times \bm H\conj) - \i\omega(\bm E \cdot \bm D\conj - \bm B \cdot \bm H\conj)}\d V.
\end{align*}
复Poynting矢量
\begin{equation}
    \bm S = \frac12\bm E \times \bm H\conj.
\end{equation}
为了分开电场项与磁场项，定义两个辅助量。无色散、无损的线性互易介质中，它们分别等于电、磁储能密度的周期平均值，
\begin{subequations}
    \begin{align}
        w_\electric &= \frac14\bm E \cdot \bm D\conj,\\
        w_\magnetic &= \frac14\bm B \cdot \bm H\conj.
    \end{align}
\end{subequations}
故简谐场的Poynting定律为
\begin{equation}
    \frac12\int_V\bm J\conj \cdot \bm E\d V + 2\i\omega\int_V(w_\electric - w_\magnetic)\d V + \oint_{\p V}\bm S \cdot \d\bm A = 0.
\end{equation}
取实部即得有功功率收支，虚部给出反应功率。若介质无损耗，$w_\electric,w_\magnetic$为实数，遂有
\begin{equation}
    \label{eq:ed-harmonic-power-balance}
    \frac12\int_V\Re(\bm J\conj \cdot \bm E)\d V
       + \oint_{\p V}\Re(\bm S) \cdot \d\bm A = 0 .
\end{equation}
此时从边界流入的平均功率全部交给显式电流。有吸收时，$2\i\omega(w_\electric - w_\magnetic)$的实部正是上节的$q_{\mathrm{abs}}$。
\paragraph{阻抗和导纳}
本节继续采用$\e^{-\i\omega t}$的时间约定。阻抗$Z$ (impedance)、电阻$R$ (resistance)、电抗$X$ (reactance)
\[
    Z = R + \i X,
\]
$X < 0$：感抗(inductance)；$X > 0$：容抗(capacitance)，因为$Z_L = -\i\omega L$、$Z_C = \i/(\omega C)$。\footnote{电气工程师的$i$已经被电流占用，故虚数单位采用j，阻抗写作$Z = R + \mathrm j X.$}

导纳$Y$ (admittance)、电导$G$ (conductance)、电纳$B$ (susceptance)
\[
    Y = G + \i B,
\]
$B < 0$：容纳(capacitive susceptance)；$B > 0$：感纳(inductive susceptance)。若改用电工学常见的$\e^{+\i\omega t}$约定，上述电抗、电纳的符号均反转。

阻抗$Z$和导纳$Y$的关系为$Y = Z\inv$。
\section{电磁场的变换性质、空间反演和时间反演}
\begin{definition}{空间反演}{spatial reflection}
    空间反演(spatial reflection)变换或称宇称(parity)变换是指
    \[
        P:\bm r_i \mapsto -\bm r_i,
    \]
    在此反演变换下，(真)矢量将同步反转(奇变换)：
    \[
        P:\bm V \mapsto -\bm V;
    \]
    而赝矢量(pseudovector)将保持不变(偶变换)：
    \[
        P:\bm A \mapsto \bm A.
    \]
\end{definition}

\begin{corollary}
    我们可以对熟知的一些矢量进行归类：
    \begin{compactitem}
    	\item 真矢量：坐标$\bm r$、速度$\bm v$、动量$\bm p$、电场$\bm E$；
    	\item 赝矢量：角动量$\bm L$、力矩$\bm N$、磁场$\bm B$。
    \end{compactitem}
    此外还有真标量(如质量$m$、电荷$q$)和赝标量(如混合积$\bm A \cdot (\bm B \times \bm C)$)之分。
\end{corollary}

\begin{remark}
    Poynting矢量是真矢量与赝矢量叉乘($\bm S = \bm E \times \bm H$)，但其仍是真矢量。
\end{remark}

\begin{definition}{时间反演}{time reversal}
    时间反演(time reversal)即
    \[
        t' = -t,
    \]
    同样有奇偶变换之分，但这与真赝矢量并无联系。
\end{definition}

将一段带电粒子的运动按相反的时间次序重放，在同一位置，电荷密度相同而电流反向。相应的源与场为
\begin{subequations}
    \label{eq:ed-time-reversal-fields}
    \begin{align}
        \rho'(\bm r,t) &= \rho(\bm r,-t),& \bm J'(\bm r,t) &= -\bm J(\bm r,-t),\\
        \bm E'(\bm r,t) &= \bm E(\bm r,-t),&\bm B'(\bm r,t) &= -\bm B(\bm r,-t).
    \end{align}
\end{subequations}
这样，速度和磁场同时反号，Lorentz力与加速度均保持原方向。真空场的能量密度也是偶的，Poynting矢量和角动量密度则反向。若原解描述向外传播的辐射，反演后的解描述向内会聚的波。选择推迟势给电磁问题指定了辐射边界条件。
